\section{Reproducibility and Artifacts}

The source repository and reviewed public evaluation artifacts are published
together \cite{openadaptartifact} at \artifacturl{}. The paper package includes
a script that reads released benchmark JSON and fails if a headline constant
drifts. Depending on the study, that released source is either a raw result or a
bounded aggregate summary. Task definitions, benchmark prose, failure
taxonomies, and selected run reports are versioned where their release boundary
permits it. Grown identity and reliability corpora, deployment-derived tuning,
target-specific recipes, and raw private rows are deliberately excluded. The
machine-check therefore establishes consistency between the paper and released
evidence; it neither claims that every underlying private row is public nor
validates the design of the underlying measurement.

Evidence carries one of four labels:
\begin{itemize}
  \item \textbf{CI-reproducible}: fixture and environment can run in automated
        CI, subject to browser and dependency availability.
  \item \textbf{Field}: real third-party target, but shared mutable state or
        credentials prevent full CI reproduction.
  \item \textbf{Fixture/fault injection}: deterministic mechanism study, not a
        population estimate.
  \item \textbf{Descriptive}: useful for discovering failures, but missing the
        repeated trials needed for comparative inference.
\end{itemize}

Table~\ref{tab:evidence} is the index: every reported study, the released
artifact the machine-check reads, its evidence label, its run count, and the
oracle that decided each run. A claim that does not appear in this table is a
mechanism description, not a measurement.

\begin{table}[t]
\centering
\caption{Evidence index. Paths are relative to the repository root; the
machine-check reads each one. ``$n$'' is runs per arm or per cell as reported in
the corresponding section.}
\label{tab:evidence}
\footnotesize
\setlength{\tabcolsep}{4pt}
\begin{tabular}{L{27mm}L{42mm}L{18mm}L{13mm}L{40mm}}
\toprule
Study & Artifact under \texttt{benchmark/} & Label & $n$ & Oracle \\
\midrule
End-to-end effect & \texttt{effect\_e2e/} \texttt{results.json} & Fixture & 9/cell & direct read-only SQLite, every table \\
Fault model (perception) & \texttt{fault\_model/} \texttt{results.json} & Fixture & 10/class & database state \\
Lending (2nd domain) & \texttt{lending\_fault\_model/} \texttt{swer\_results.json} & Fixture & 3/task & separate read-only ledger \\
AP invoice & \texttt{ap\_invoice/} \texttt{results.json} & Fixture & 3/cell & two-database delta audit \\
O2C reconciliation & \texttt{o2c\_recon/} \texttt{results.json} & Fixture & 3/cell & two-database delta audit \\
OpenEMR comparison & \texttt{openemr/} \texttt{results.json} & Field & 20 / 10 & task-level success check \\
MockMed comparison & \texttt{results.json} & CI-repro. & 100 / 20 & task-level success check \\
Public-web breadth & \texttt{reliability/} \texttt{summary.json} & Descriptive & 1/app & external goal oracle \\
Identity ladder & \texttt{identity\_ladder/} \texttt{identity\_ladder.json} & Fixture & 14--42 & intended-entity ground truth \\
Windows UIA & \texttt{windows\_uia/} \texttt{results.json} & Field & 3 & independent SQLite query \\
Native macOS & \texttt{macos\_native/} \texttt{*.adjudication.json} & Field & 3 & exact file bytes \\
Network RDP & \texttt{rdp/} \texttt{*.sanitized.json} & Field & 3 & guest-tools file readback \\
Citrix stand-in & \texttt{citrix\_ica\_hdx/} \texttt{results.json} & Fixture & 29 trials & in-process out-of-band store \\
\bottomrule
\end{tabular}
\end{table}

Reproducibility is uneven across these labels, and the unevenness is
structural rather than incidental. The fault-injection, end-to-end effect,
lending, and multi-application studies are localhost-only, deterministic, and
make no model calls, so they re-run from a clean checkout and reproduce exactly.
The comparative study does not: it depends on a hosted model identified by an
exact name, tool version, and dated beta header, none of which the authors
control or can guarantee will remain available or unchanged, so its latency and
cost figures are pinned to a date rather than reproducible on demand. The
OpenEMR arm additionally runs against a shared, mutable, public demonstration
instance and is not reproducible by design. The desktop and remote-substrate
qualifications depend on specific virtual machines, snapshots, and hosts. We
report these as field or descriptive evidence rather than presenting the whole
suite as uniformly re-runnable.

Before submission, the release candidate must re-run promoted experiments,
archive its exact commit and environment, record task-level oracles, and check
that no application artifact contains sensitive or prohibited data. The
repository's \texttt{paper/ARTIFACT\_CHECKLIST.md} is the submission gate, and
records the exact commit, release tag, and published artifact digests that this
version of the paper describes.

\paragraph{Data availability and ethics.}
Every released system-of-record fixture is synthetic: the MockMed and MockLoan
fixtures, their injected-fault suites, and the benchmark task packs contain no
real patient, borrower, or customer data. The public-web study observed public,
unauthenticated applications, but only its bounded aggregate is released; its
raw per-target rows and curated target recipes remain private. This paper
release contains no customer-derived hardening corpus. Such grown corpora and
raw private evaluation rows remain outside the public repository; when an
approved sanitized derivative crosses a trust boundary, the pipeline applies
type-specific transformations, rescans the result, and binds human approval to
its exact hash as described in
Section~\ref{sec:governance}. Human review reduces but does not prove
de-identification, so sanitization remains a defense-in-depth control rather
than a guarantee. Repository-only third-party reference environments are also
excluded from the permissively licensed wheel and source distribution; the
release gate inspects the built archives rather than inferring their contents
from the source tree.
